\section{Interpolation with separated weights}
\label{geom:section}

We prove the interpolation statement used in the determinant argument. Its
essential uniformity is that the weights can be chosen before the interpolation
points: each coordinate weight has a lower threshold depending on the earlier
weights, but not on the locations of the points. The degree at which
interpolation becomes surjective is allowed to depend on all the fixed data.

All multi-indices below have nonnegative integer entries. For a positive
rational vector $W=(w_0,\ldots,w_m)$, the $W$-degree of the
monomial $Y^hX_1^{\alpha_1}\cdots X_m^{\alpha_m}$ is
$w_0h+\sum_iw_i\alpha_i$. Write $\mathcal P_W(H)$ for the complex vector
space of polynomials of $W$-degree at most $H$. For a positive rational vector
$V=(v_0,\ldots,v_m)$, let $\mathcal J_V(H)$ be the quotient of
$\CC[[t,u_1,\ldots,u_m]]$ by the ideal of series whose monomials all have
$V$-weight at least $H$. Thus a vector in $\mathcal J_V(H)$ is a packet of
coefficients indexed by
\[
 v_0s+\sum_{i=1}^m v_i\beta_i<H.
\]
The strict inequality in this definition will be used throughout.

\begin{theorem}[Separated-weight interpolation]
\label{geom:interpolation}
Fix integers $m,K\ge1$ and positive rational numbers $w_0,v_0,\theta$ such
that $0<\theta<1$ and
\begin{equation}
\label{geom:volume-hypotheses}
 K\theta^m<1,
 \qquad K\frac{w_0}{v_0}\theta^m<1.
\end{equation}
There are successive lower thresholds for positive rational numbers
$w_1,\ldots,w_m$, depending only on these fixed data and the previously
chosen weights, with the following property. Put
\[
 W=(w_0,w_1,\ldots,w_m),\qquad
 V=(v_0,w_1/\theta,\ldots,w_m/\theta).
\]
For any points
\[
 a_j=(1,c_{j1},\ldots,c_{jm})\in\CC^{m+1},\qquad 0\le j<K,
\]
whose coordinates $c_{0i},\ldots,c_{K-1,i}$ are pairwise distinct for each
$i$, the map
\begin{equation}
\label{geom:jet-map}
 \begin{aligned}
 \mathcal P_W(H)&\longrightarrow\bigoplus_{j=0}^{K-1}\mathcal J_V(H),\\
 P&\longmapsto
 \left(P\bigl(1+t,(c_{ji}+u_i+\log(1+t))_{i=1}^m\bigr)\right)_{j=0}^{K-1}.
 \end{aligned}
\end{equation}
is surjective for all sufficiently large, sufficiently divisible integers $H$.
The divisibility requirement can be chosen using only the weights. The
remaining lower threshold for $H$ may depend on the points.
\end{theorem}

We first establish a numerical inequality on algebraic curves, and then use
it to obtain all the jets in \eqref{geom:jet-map}. The local estimate driving
the curve argument is recorded separately.

\subsection{A local multiplicity comparison}

A collection of vector fields is called a \emph{normal basis} to a smooth
subvariety at a point if its images form a basis of the ambient tangent space
modulo the tangent space of the subvariety. For an irreducible subvariety,
this terminology will refer to a general smooth point. There are only
finitely many possible index sets in the applications below, so all their
generic rank conditions can be checked on one common open subset.

\begin{lemma}[Weighted transverse multiplicity]
\label{geom:multiplicity}
Give the coordinates $x_1,\ldots,x_d$ positive rational degree weights
$\rho_1,\ldots,\rho_d$. On an open subset of $\CC^d$, let
$D_1,\ldots,D_d$ be commuting analytic vector fields forming a frame, and
assign them positive rational costs $\kappa_1,\ldots,\kappa_d$. Fix
$\varepsilon>0$.

Let $N>0$. Suppose that polynomials $f_1,\ldots,f_r$ have weighted degree at most $N$,
that $Z$ is an irreducible component of their common zero set of codimension
$k\ge1$, and that
\begin{equation}
\label{geom:persistent-derivatives}
 D^\gamma f_\ell\big|_Z=0
 \quad\text{whenever}\quad
 \sum_b\kappa_b\gamma_b<\varepsilon N.
\end{equation}
For any coordinate normal basis $A$ and any $D$-field normal basis $B$,
each of cardinality $k$, one has
\begin{equation}
\label{geom:multiplicity-comparison}
 \prod_{a\in A}\rho_a
 \le 2k!\,\varepsilon^{-k}\prod_{b\in B}\kappa_b
\end{equation}
The bound is uniform in the equations, the component $Z$, and the
analytic frame; no lower threshold for $N$ is required.
\end{lemma}

\begin{proof}
Choose a smooth point $x$ of $Z$ where both normal-basis conditions hold and
which lies on no other irreducible component of the common zero set. Flowing
from $Z$ along the fields indexed by $B$ gives local analytic coordinates
$(z,b)$, where $z$ parametrizes $Z$ and $b=(b_1,\ldots,b_k)$ consists of
transverse flow parameters. Commutativity implies that the chosen fields
are the derivatives in the $b$-directions. Thus
\eqref{geom:persistent-derivatives} says that in the expansion
\[
 f_\ell(z,b)=\sum_{\alpha\in\NN^k} f_{\ell,\alpha}(z)b^\alpha,
\]
every coefficient function of $B$-weight less than $\varepsilon N$ vanishes
identically on the chosen neighborhood in $Z$.

Let $S$ be the affine coordinate slice through $x$ obtained by fixing the
coordinates outside $A$. It is transverse to $Z$, and the zero set of the
restricted equations is isolated at $x$. The flow parameters $b$ give
regular analytic coordinates on $S$; the coordinates $z$ become analytic
functions $z(b)$. Substitution cannot generate a missing term of smaller
$B$-weight. Indeed, each surviving summand already contains a monomial
$b^\alpha$ of weight at least $\varepsilon N$, and its analytic coefficient
has only nonnegative powers of $b$. Consequently the ideal of the
restricted equations is contained in the monomial ideal of weight at least
$\varepsilon N$. Its quotient therefore has length at least
\begin{equation}
\label{geom:local-length-lower}
 \#\left\{\alpha\in\NN^k:
       \sum_{b\in B}\kappa_b\alpha_b<\varepsilon N\right\}
 \ge \frac{(\varepsilon N)^k}{2k!\prod_{b\in B}\kappa_b}
\end{equation}
In fact the factor $1/2$ can be omitted: the simplex
$\{x\in\RR_{\ge0}^k:\sum_{b\in B}\kappa_bx_b<\varepsilon N\}$
is covered by the disjoint half-open unit cubes $\alpha+[0,1)^k$
indexed by the admissible multi-indices $\alpha=\lfloor x\rfloor$. Its volume is
$(\varepsilon N)^k/(k!\prod_{b\in B}\kappa_b)$. Thus the displayed
bound holds for every $N>0$. Analytic local coordinates preserve finite
length, so no radius or coefficient bound for the flow map is needed.

In the regular local ring of the $k$-dimensional slice, the ideal generated
by the restricted equations is primary to the maximal ideal. Over the
infinite field $\CC$, $k$ sufficiently general constant linear
combinations of these equations still have an isolated common zero at
$x$: successive generic combinations avoid the finitely many
positive-dimensional components that remain. Their ideal is contained
in the original ideal. Their quotient length is therefore an upper bound
for the length in \eqref{geom:local-length-lower}.

Each of the $k$ combinations has weighted degree at most $N$ in the
variables indexed by $A$. We claim that their local intersection length
is at most
\begin{equation}
\label{geom:weighted-bezout}
 \frac{N^k}{\prod_{a\in A}\rho_a}.
\end{equation}
Choose an integer $M_0$ clearing the denominators of the $\rho_a$, and
make the substitutions
\[
 x_a=\lambda_a+z_a^{M_0\rho_a}\qquad(a\in A).
\]
Choose the constants $\lambda_a$ so that no preimage of the given point
is ramified. There are $\prod_{a\in A}(M_0\rho_a)$ such preimages, each
with the same local intersection length. The substituted polynomials
have ordinary total degree at most $M_0N$. The ordinary B\'ezout bound
for isolated complete-intersection zeros, with their local
multiplicities, bounds the sum of these lengths by $(M_0N)^k$;
see \cite[Corollary~VIII.3, p.~144]{mondal2021}.
The multiplicities there are local quotient lengths, unchanged by
completion \cite[Section~IV.4.2]{mondal2021}.
Dividing gives \eqref{geom:weighted-bezout}.
This application concerns isolated zeros only and permits other
components elsewhere.

Combining \eqref{geom:local-length-lower} and
\eqref{geom:weighted-bezout} proves the result. Neither the
polynomials nor the possibly varying component enters the constant in
\eqref{geom:multiplicity-comparison}.
\end{proof}

\subsection{The curve inequality}

Choose a positive rational number $\sigma$ sufficiently small that
\begin{equation}
\label{geom:sigma}
 \begin{split}
 (1+3\sigma)^{m+1}K(w_0/v_0)\theta^m&<1,\\
 (1+3\sigma)^mK\theta^m&<1,\\
 (1+\sigma)\theta&<1.
 \end{split}
\end{equation}
This choice is possible by \eqref{geom:volume-hypotheses}. Set $v_i=w_i/\theta$
for $i>0$. Fix $L_0=m+2$ and put
\begin{equation}
\label{geom:constant}
 C_* =\max_{1\le k\le m}2k!(L_0/\sigma)^k.
\end{equation}
We choose the positive weights successively so that for every pair of
subsets $A,B\subseteq\{0,\ldots,m\}$ of the same cardinality between
$1$ and $m$,
\begin{equation}
\label{geom:separation}
 \prod_{a\in A}w_a>C_*\prod_{b\in B}v_b
\end{equation}
whenever the largest positive index in their symmetric difference belongs
to $A$.

Here and below a condition involving that largest index is imposed only
when a positive index does differ. The successive choice is possible:
if this largest index is $i$, the ratio in \eqref{geom:separation} is
$w_i$ times a positive expression in the already chosen lower weights
and the fixed zeroth weights. Every higher index occurs in both sets or
in neither; in the former case its ratio is $w_j/v_j=\theta$. Hence the
expression is independent of all higher weights. There are finitely many
pairs of sets. A sufficiently large choice of $w_i$ enforces all
conditions whose largest differing index is $i$, and subsequent choices
do not disturb them. This construction involves no interpolation point.

Let $C$ be an irreducible algebraic curve whose generic point lies in the
affine coordinate chart. Such a curve has a smooth projective model over
$\CC$; see \cite[Theorem~53.2.6 and Lemma~53.2.8, Tag~0BXX]{stacks}.
On that model define
\begin{equation}
\label{geom:curve-degree}
 \deg_W C=
 \sum_P\max\left\{0,-\frac{\operatorname{ord}_P Y}{w_0},
                  -\frac{\operatorname{ord}_P X_1}{w_1},\ldots,
                  -\frac{\operatorname{ord}_P X_m}{w_m}\right\}.
\end{equation}
The order of an identically zero coordinate is understood as $+\infty$;
it contributes no pole. At a branch $P$ through $a_j$, set
\begin{equation}
\label{geom:branch-contact}
 h_P=\min\left\{
 \frac{\operatorname{ord}_P(Y-1)}{v_0},
 \frac{\operatorname{ord}_P(X_1-c_{j1}-\log Y)}{v_1},\ldots,
 \frac{\operatorname{ord}_P(X_m-c_{jm}-\log Y)}{v_m}
 \right\}.
\end{equation}
The logarithm here is its formal expansion at $Y=1$. All entries have
positive order, possibly infinite, and their minimum is finite on a
nonconstant curve.

\begin{proposition}[Curve inequality]
\label{geom:curve-inequality}
With the weights chosen above, every such curve satisfies
\begin{equation}
\label{geom:curve-bound}
 \deg_W C\ge(1+\sigma)\sum_{P\mapsto a_j}h_P,
\end{equation}
where the sum includes all branches through the prescribed points.
\end{proposition}

\begin{proof}
There is nothing to prove if the curve misses all the points. Suppose
instead that \eqref{geom:curve-bound} fails. For every sufficiently large
auxiliary integer $N$, linear algebra gives a nonzero polynomial $F_N$
of $W$-degree at most $N$ whose expansions at all the points have
$V$-order at least $(1+3\sigma)N$. Indeed, the number of its coefficients
has leading term
\[
 \frac{N^{m+1}}{(m+1)!\prod_{a=0}^m w_a},
\]
whereas the number of conditions has leading term
\[
 \frac{K(1+3\sigma)^{m+1}N^{m+1}}
      {(m+1)!\prod_{a=0}^m v_a}.
\]
Their comparison is the first inequality in \eqref{geom:sigma}. This
argument requires no independence among the conditions.

On the open set $Y\ne0$ consider the commuting frame
\begin{equation}
\label{geom:frame}
 D_0=Y\partial_Y+\sum_{i=1}^m\partial_{X_i},\qquad
 D_i=\partial_{X_i}\quad(1\le i\le m).
\end{equation}
It preserves polynomial degree at most $N$. In the formal coordinates of
\eqref{geom:jet-map},
\[
 D_0=(1+t)\partial_t,\qquad D_i=\partial_{u_i}.
\]
Thus differentiation by $D_a$ decreases weighted order by at most $v_a$.
Every $D^\gamma F_N$ with $\sum_av_a\gamma_a\le\sigma N$ retains
order at least $(1+2\sigma)N$ at each point. On a branch $P$, its order
is consequently at least $(1+2\sigma)N h_P$.

A polynomial of $W$-degree at most $N$ restricts to a rational function
on the complete model of $C$ with total pole degree at most
$N\deg_W C$. Since we assumed that \eqref{geom:curve-bound} fails, the
just obtained total zero order exceeds this pole bound. The ordinary
divisor equality for a nonzero rational function therefore forces
\begin{equation}
\label{geom:derivatives-on-curve}
 D^\gamma F_N\big|_C=0
 \qquad\left(\sum_av_a\gamma_a\le\sigma N\right).
\end{equation}
All the assertions about branch orders hold for formal power series:
each monomial has order at least its weight times $h_P$, and
cancellation can only increase the order.

We have converted excess contact into vanishing of many derivatives on
$C$. We next find a component on which that vanishing persists, as in
the derivative-ideal method of zero estimates
\cite[Section~5]{philippon1986}, and compare all its normal bases.
Put $\delta=\sigma N/L_0$. For $0\le r\le L_0$, let $V_r$ be the common
zero set in $Y\ne0$ of the derivatives of cost at most $r\delta$.
Let $C^\circ$ be the part of $C$ in this affine open set. The $V_r$
form a decreasing chain, and \eqref{geom:derivatives-on-curve} puts
$C^\circ$ in $V_{L_0}$. Choose an irreducible component of $V_{L_0}$ containing
$C^\circ$, and extend it backwards to nested irreducible components of the
earlier $V_r$. Every selected component has dimension between $1$ and
$m$: it contains a curve, and it lies in the proper zero set of
$F_N$. Since there are $L_0+1=m+3$ levels, two adjacent selected
components coincide. Denote this persistent component by $Z$, and
write $k=\operatorname{codim}Z$.

Take the derivatives of cost at most $r\delta$ as the equations of the
earlier of these two levels. Any further derivative of cost less than
$\delta$ vanishes along $Z$, because the fields commute and $Z$ is
also a component at the next level. Lemma~\ref{geom:multiplicity}, with
$\varepsilon=\sigma/L_0$, gives
\begin{equation}
\label{geom:all-bases}
 \prod_{a\in A}w_a\le C_*\prod_{b\in B}v_b
\end{equation}
for every coordinate normal basis $A$ and every frame normal basis $B$.
The equations and $Z$ may vary with $N$; the uniformity proved in
Lemma~\ref{geom:multiplicity} is exactly what permits this use.

We show that $Y$ is constant on $Z$, working on the common smooth open
subset where all the normal-basis conditions hold. Suppose first that
some positive coordinate direction is nontangent to $Z$, and choose the
largest such index $i$. Every larger positive index occurs in neither
kind of normal basis, since $D_j=\partial_{X_j}$ for $j>0$.
Some coordinate normal basis $A$ contains $i$. If a frame normal basis
$B$ omitted $i$, the largest positive index where these sets differ
would be $i$. Then
\eqref{geom:separation} would contradict \eqref{geom:all-bases}.
Thus every frame normal basis contains $D_i$.

The other frame vectors consequently fail to span the normal quotient.
Their ambient span is the hyperplane
\[
 \ker(dX_i-dY/Y).
\]
For an ambient hyperplane $Q$ and a tangent space $T$, the image of
$Q$ in the quotient by $T$ fails to be surjective precisely when
$T\subseteq Q$. We obtain
\begin{equation}
\label{geom:log-differential}
 dX_i=dY/Y\quad\text{on }Z.
\end{equation}
This identity forces $Y$ to be constant on the algebraic variety $Z$.
Indeed, if $Y$ were nonconstant, choose a smooth point where $dY\ne0$
and intersect with general algebraic hyperplanes through that point
until a curve remains on which $Y$ is nonconstant. On the smooth
complete model of this curve, \eqref{geom:log-differential} is an
identity of meromorphic differentials. A nonconstant meromorphic
function $Y$ has a zero or pole, where $dY/Y$ has nonzero residue
$\operatorname{ord}_P Y$. An exact differential $dX_i$ has zero
residue at every point. This is a contradiction. It follows also
that $X_i$ is constant, since $dX_i=0$ in characteristic zero.

If instead every positive coordinate direction is tangent, the $m$
independent vectors $\partial_{X_i}$ all lie in the tangent space of
$Z$. Properness forces that tangent space to be their span, and
$dY=0$ on $Z$ again. We conclude in all cases that $Y$ is constant
on $Z$. Since $C^\circ\subseteq Z$ meets a prescribed point with
$Y=1$, we have $Y=1$ on $C^\circ$, and hence on $C$ by density.

It remains to rule out a violating curve in that fiber. We work with
$C^\circ$ in the affine fiber $Y=1$, while computing orders and degrees
on the complete model of $C$ as before. In this fiber, the degree weights
are $w_i$, the order weights are $v_i=w_i/\theta$,
and $h_P=\min_i\operatorname{ord}_P(X_i-c_{ji})/v_i$.
Choose a new auxiliary polynomial in these $m$ variables. The ratio of
the number of vanishing conditions to the number of coefficients now
tends to $K(1+3\sigma)^m\theta^m<1$, by the second inequality
in \eqref{geom:sigma}. Thus, for sufficiently large $N$, there is a
nonzero polynomial $G_N(X_1,\ldots,X_m)$ of weighted degree at most $N$
and order at least $(1+3\sigma)N$ at every center.

An ordinary partial derivative of cost at most $\sigma N$ still has
order at least $(1+2\sigma)N$ at every center and degree at most $N$.
The same zero--pole comparison on $C$ therefore gives
\[
 \partial_X^\gamma G_N\big|_C=0
 \qquad\left(\sum_{i=1}^m v_i\gamma_i\le\sigma N\right).
\]
If $m=1$, this is already impossible: a nonzero polynomial on the affine
line cannot vanish on a curve. Suppose henceforth that $m\ge2$. For
$0\le r\le L_0$, let $U_r$ be the common zero set in this fiber of the
partial derivatives of $G_N$ of cost at most $r\delta$, with
$\delta=\sigma N/L_0$ as before. The curve $C^\circ$ lies in $U_{L_0}$. Choose
nested irreducible components containing $C^\circ$ backwards through this
chain. Their dimensions lie between $1$ and $m-1$, since $G_N$ is
nonzero, so two adjacent components coincide in a component $Z'$.
Its codimension in the fiber satisfies $1\le k\le m-1$.

For the equations at the earlier of these two levels, every further
partial derivative of cost less than $\delta$ vanishes on $Z'$.
Lemma~\ref{geom:multiplicity}, again with $\varepsilon=\sigma/L_0$,
therefore gives \eqref{geom:all-bases} for $Z'$, with the same $C_*$
and with normal-basis index sets contained in $\{1,\ldots,m\}$.

Here the coordinate and frame bases are the same. They must have a
unique index set: two different basis sets could be ordered so that
the largest differing index belongs to $A$, contradicting
\eqref{geom:separation}. The codimension is positive, so this unique
normal basis has some index $i$. The other coordinate directions
fail to span the normal quotient; hence the tangent space is
contained in $\ker dX_i$, and $X_i$ is constant on $Z'$.
Coordinatewise distinctness shows that $Z'$ meets at most one prescribed
point. Every prescribed point of $C$ belongs to $C^\circ\subseteq Z'$,
so the same holds for $C$.

Let that point be $a_j$. Choose a coordinate $X_l$ nonconstant on
$C$. The nonzero rational function $X_l-c_{jl}$ has total pole degree
at most $w_l\deg_W C$, and at each branch through $a_j$ it has order
at least $(w_l/\theta)h_P$. Therefore
\[
 \deg_W C\ge\theta^{-1}\sum_{P\mapsto a_j}h_P
             >(1+\sigma)\sum_{P\mapsto a_j}h_P,
\]
by the last inequality in \eqref{geom:sigma}. This final contradiction
proves \eqref{geom:curve-bound}.
\end{proof}

\subsection{From the curve inequality to all coefficient packets}

We complete the proof of Theorem~\ref{geom:interpolation}. The argument uses
ordinary projective schemes and the ordinary blow-up; smoothness or
normality of the compactification will not be needed. Curve positivity
and asymptotic jet generation are closely related in the Seshadri-constant
approach of \cite[Section~6]{demailly1992}. Here we supply the passage
for our weighted ideals at several points, including the eventual
ordinary-power comparison required on a possibly singular scheme.

Choose an integer $R$ sufficiently large that $R/w_a$ and $R/v_a$ are
positive integers for every index. Embed affine space by all monomials
of $W$-degree at most $R$, including the constant monomial, and let
$X$ be the projective closure of the image. Since these monomials
include each coordinate, the chart of the constant monomial is
precisely the original affine space. Let $L=\mathcal O_X(1)$, a very
ample line bundle.

For any curve meeting this chart,
\begin{equation}
\label{geom:projective-degree}
 \deg(L|_C)=R\deg_W C.
\end{equation}
To see this on its smooth complete model, at each point clear the
largest pole of the homogeneous monomial coordinates. Every monomial
has pole order at most $R$ times the maximum in
\eqref{geom:curve-degree}. A pure power of one coordinate attains
that bound whenever the maximum is positive, and the constant
monomial attains it when the maximum is zero. The resulting sections
have no common zero, proving \eqref{geom:projective-degree}.

To make the local data algebraic, replace $\log(1+t)$ in each
coordinate by a fixed Taylor polynomial $G_i(t)$ whose omitted
terms have $t$-weight strictly greater than $v_i$. Then
\[
 t=Y-1,\qquad u_i'=X_i-c_{ji}-G_i(t)
\]
are regular local coordinates at $a_j$. This replacement preserves
$h_P$. Indeed, if $t$ is not identically zero on a branch, each
omitted term has order divided by $v_i$ strictly larger than
$\operatorname{ord}_P(t)/v_0$; if $t$ is identically zero, its tail
vanishes. It also preserves interpolation surjectivity for every
threshold: the formal coordinate change fixes $t$ and changes
$u_i$ only by terms of weight greater than $v_i$. It and its inverse
preserve the weighted filtration and induce inverse triangular maps
on each finite coefficient packet.

At $a_j$ form the finite-colength ideal
\[
 I_j=\left(t^{R/v_0},(u_1')^{R/v_1},\ldots,(u_m')^{R/v_m}\right).
\]
The points are distinct. These local ideals, together with the unit
ideal away from the points, define a coherent ideal sheaf $I$ on $X$.
For a branch through a center,
\begin{equation}
\label{geom:ideal-valuation}
 \operatorname{ord}_P I
 =\min_a\left\{(R/v_a)\operatorname{ord}_P z_a\right\}
 =Rh_P,
 \qquad z_0=t,\quad z_i=u_i'.
\end{equation}
Let $p:X'\to X$ be the ordinary blow-up of $I$, namely the relative
Proj of its Rees algebra, and write
\[
 I\mathcal O_{X'}=\mathcal O_{X'}(-E).
\]
The scheme $X$ is integral and Noetherian, $I$ is a coherent nonzero
ideal, and $E$ is an effective Cartier divisor. The tautological
bundle $\mathcal O_{X'}(-E)$ is relatively ample; these are the
standard blow-up properties in
\cite[Tags~02OS, 02ND, and~02NS]{stacks}.

Write $A=p^*L$ in additive divisor notation. We claim that
\begin{equation}
\label{geom:nef-class}
 A-(1+\sigma)E\quad\text{is nef}.
\end{equation}
For a noncontracted integral curve $\Gamma$ upstairs whose image
meets the affine chart, the map to its image is birational, since
$p$ is an isomorphism away from the finite center. On the complete
normalizations, \eqref{geom:projective-degree} and
\eqref{geom:ideal-valuation} give
\[
 A\cdot\Gamma=R\deg_W p(\Gamma),\qquad
 E\cdot\Gamma=R\sum_{P\mapsto a_j}h_P.
\]
Proposition~\ref{geom:curve-inequality} gives the desired nonnegative
degree. A curve whose image misses the affine chart avoids the
exceptional divisor. A contracted curve has $A$-degree zero and
positive degree for $-E$ by relative ampleness. These cases account
for every integral curve, proving \eqref{geom:nef-class}.

For sufficiently large $a>1$, the divisor $aA-E$ is ample. One
can see this directly: choose $b$ so that $I\otimes L^b$ is globally
generated. The Rees construction embeds $X'$ in a product
$X\times\mathbb P^r$, with the second hyperplane bundle restricting
to $bA-E$. Tensoring with the first very ample bundle makes
$(b+1)A-E$ ample. Now the identity
\begin{equation}
\label{geom:ample-combination}
 A-E=
 \frac{a-1}{a(1+\sigma)-1}\bigl(A-(1+\sigma)E\bigr)
 +\frac{\sigma}{a(1+\sigma)-1}(aA-E)
\end{equation}
expresses $A-E$ as a positive multiple of a nef class plus a
positive multiple of an ample class. It is therefore ample;
see \cite[Corollary~1.4.10]{lazarsfeld}. This theorem applies to
projective schemes and does not require a smooth blow-up.

For all sufficiently large $n$, the graded-piece and relative
vanishing theorems for the ordinary Rees algebra give
\begin{equation}
\label{geom:rees-pushdown}
 p_*\mathcal O_{X'}(-nE)=I^n,\qquad
 R^q p_*\mathcal O_{X'}(-nE)=0\quad(q>0).
\end{equation}
Indeed, on any affine open of the Noetherian base the Rees algebra
is a Noetherian graded algebra generated in degree one. Its
sufficiently high graded pieces recover the sections of its
associated sheaf on Proj, and the higher cohomology of those twists
vanishes. A finite affine cover gives one common threshold;
see \cite[Tags~0AG6 and~0AG7]{stacks}. In particular,
\eqref{geom:rees-pushdown} uses ordinary powers $I^n$, not their
integral closures, and makes no assertion about small $n$.

By the projection formula, Leray, and Serre vanishing for the ample
line bundle $A-E$,
\begin{equation}
\label{geom:serre-vanishing}
 H^1(X,I^n\otimes L^n)
 =H^1(X',\mathcal O_{X'}(n(A-E)))=0
\end{equation}
for all sufficiently large $n$; see \cite[Tag~0B5U]{stacks}.
The exact sequence for $I^n$ consequently gives a surjection from
$H^0(X,L^n)$ onto all its local coefficient data modulo $I^n$.
The constant projective coordinate trivializes $L$ on the affine
chart. Finite-colength local quotients are unchanged by completion,
so these algebraic classes are also the formal coefficient data used
in the theorem. In these trivializations, every element of $I^n$ has
weighted order at least $nR$, because every displayed generator
of $I$ has weight $R$. Thus
\begin{equation}
\label{geom:ideal-inclusion}
 I^n\subseteq\{\text{local series of weight at least }nR\}.
\end{equation}
Surjectivity modulo $I^n$ therefore implies surjectivity onto the
smaller quotient consisting of the desired packets of weight less
than $nR$. Equality in \eqref{geom:ideal-inclusion} is not needed.

Finally, every section of $L^n$ for sufficiently large $n$ is
supplied by a homogeneous degree-$n$ polynomial in the projective
embedding coordinates. This follows from Serre vanishing applied
to the homogeneous ideal sheaf of $X$ in its projective space.
Such a polynomial restricts on the affine chart to one of
$W$-degree at most $nR$. Taking $H=nR$ and undoing the filtered
Taylor-polynomial coordinate change proves the surjectivity in
\eqref{geom:jet-map}, and completes the proof of
Theorem~\ref{geom:interpolation}.
